% ============================================================================
%  Manual de uso - S&A Tool Hub (SyA-hub)
%  Compilar con:  pdflatex Manual_SyA_Hub.tex   (ejecutar dos veces para el índice)
% ============================================================================
\documentclass[11pt,a4paper]{article}

\usepackage[utf8]{inputenc}
\usepackage[T1]{fontenc}
\usepackage[spanish,es-noshorthands]{babel}
\usepackage{lmodern}
\usepackage[margin=2.3cm]{geometry}
\usepackage[table]{xcolor}
\usepackage{listings}
\usepackage{booktabs}
\usepackage{tabularx}
\usepackage{enumitem}
\usepackage{fancyhdr}
\usepackage{titlesec}
\usepackage[most]{tcolorbox}
\usepackage[hidelinks]{hyperref}

% ---------------------------------------------------------------- Colores ---
\definecolor{sablue}{HTML}{002395}
\definecolor{sanavy}{HTML}{001360}
\definecolor{salight}{HTML}{EEF3FF}
\definecolor{codebg}{HTML}{F6F8FA}
\definecolor{warnbg}{HTML}{FFF7E6}
\definecolor{warnfr}{HTML}{D97706}

\titleformat{\section}{\Large\bfseries\color{sanavy}}{\thesection}{0.8em}{}
\titleformat{\subsection}{\large\bfseries\color{sablue}}{\thesubsection}{0.7em}{}

\pagestyle{fancy}
\fancyhf{}
\fancyhead[L]{\small\color{sablue}\textbf{S\&A Tool Hub}}
\fancyhead[R]{\small Manual de uso}
\fancyfoot[C]{\small\thepage}
\renewcommand{\headrulewidth}{0.4pt}

\lstset{
  basicstyle=\ttfamily\small,
  backgroundcolor=\color{codebg},
  frame=single, rulecolor=\color{gray!40},
  breaklines=true, columns=fullflexible,
  keepspaces=true, showstringspaces=false,
  xleftmargin=4pt, xrightmargin=4pt,
  literate={á}{{\'a}}1 {é}{{\'e}}1 {í}{{\'i}}1 {ó}{{\'o}}1 {ú}{{\'u}}1 {ñ}{{\~n}}1
}

\newtcolorbox{nota}[1][Nota]{colback=salight,colframe=sablue,fonttitle=\bfseries,title=#1,arc=2mm}
\newtcolorbox{aviso}[1][Advertencia]{colback=warnbg,colframe=warnfr,fonttitle=\bfseries,title=#1,arc=2mm}

\newcommand{\ruta}[1]{\texttt{#1}}

% ============================================================================
\begin{document}

% ---------------------------------------------------------------- Portada ---
\begin{titlepage}
  \centering
  \vspace*{3cm}
  {\fontsize{40}{48}\selectfont\bfseries\color{sablue} S\&A}\\[2pt]
  {\large\bfseries\color{sablue} CHILE}\\[1.5cm]
  {\Huge\bfseries\color{sanavy} S\&A Tool Hub}\\[0.4cm]
  {\Large Centro de Herramientas de Análisis}\\[1.8cm]
  {\large Manual de instalación, uso y características}\\[0.3cm]
  {\large Control de acceso por roles (RBAC) y credenciales}\\
  \vfill
  {\small Versión del documento: 1.0 \quad---\quad Noviembre 2026}
\end{titlepage}

% ------------------------------------------------- Aviso de personalización ---
\begin{tcolorbox}[colback=salight,colframe=sablue,fonttitle=\bfseries\large,
  title={Este software es completamente modificable},arc=2mm]
S\&A Tool Hub se entrega como \textbf{código fuente abierto para el equipo}: puede
modificarse, ampliarse o recortarse \textbf{a gusto}. Textos, colores, logotipos,
módulos, umbrales, usuarios, roles, endpoints de la API y el propio flujo de navegación
están en archivos TypeScript/React legibles y documentados en este manual.

Nada de lo descrito aquí es una restricción: es la configuración \emph{actual} del
proyecto, pensada como punto de partida. Las secciones~\ref{sec:estructura},
\ref{sec:rbac} y~\ref{sec:credenciales} indican exactamente qué archivo tocar para
cada cambio habitual.
\end{tcolorbox}

\vspace{0.6cm}
\tableofcontents
\newpage

% ============================================================================
\section{Descripción general}

S\&A Tool Hub es una aplicación web (SPA en React servida por Express) que reúne
herramientas internas de análisis y un portal de informes para clientes. Actualmente
ofrece:

\begin{itemize}[leftmargin=*]
  \item \textbf{Intérprete Perfmon SQL}: carga logs de Perfmon de SQL Server en formato
        CSV, genera gráficos interactivos por contador y detecta violaciones de umbral.
        Dispone de dos módulos de contexto: \textbf{SII} y \textbf{TGR}.
  \item \textbf{Exportación de reportes} del análisis en CSV y JSON, y descarga de cada
        gráfico como imagen PNG.
  \item \textbf{Repositorio local} de archivos analizados (en el \texttt{localStorage}
        del navegador): seleccionar, renombrar y eliminar análisis previos.
  \item \textbf{Historial de acciones} con las operaciones realizadas en la sesión.
  \item \textbf{Notas y tareas} con estados \emph{Pendiente}, \emph{En Proceso} y \emph{Listo}.
  \item \textbf{Portal de clientes}: cada cliente ve únicamente los informes de su entidad,
        con búsqueda, filtro por tipo y descarga.
  \item \textbf{Panel de administración}: alta, edición, cambio de rol y baja de usuarios.
\end{itemize}

\subsection{Tecnologías}

\begin{center}
\rowcolors{2}{salight}{white}
\begin{tabularx}{\linewidth}{lX}
\toprule
\textbf{Componente} & \textbf{Tecnología} \\
\midrule
Frontend        & React 19, TypeScript, Tailwind CSS 4, iconos \texttt{lucide-react} \\
Gráficos        & Chart.js 4 + \texttt{react-chartjs-2} \\
Servidor / API  & Node.js + Express 4 (archivo \ruta{server.ts}) \\
Empaquetado     & Vite 6 (frontend) y esbuild (servidor) \\
Protección      & \texttt{vite-plugin-javascript-obfuscator}: ofusca el JavaScript del build \\
\bottomrule
\end{tabularx}
\end{center}

\begin{nota}
Las versiones indicadas son las declaradas en \ruta{package.json} al momento de
redactar este manual. Pueden cambiar si se actualizan las dependencias; verifique
siempre ese archivo como fuente de verdad.
\end{nota}

% ============================================================================
\section{Instalación y puesta en marcha}

\subsection{Requisitos previos}
\begin{itemize}[leftmargin=*]
  \item \textbf{Node.js} 20 o superior (el proyecto se probó con Node.js 26) y \textbf{npm}.
        Descarga oficial: \url{https://nodejs.org}.
  \item \textbf{Git} para clonar el repositorio.
  \item Un navegador moderno (Chrome, Edge o Firefox).
\end{itemize}

\subsection{Obtener el código}
\begin{lstlisting}[language=bash]
git clone https://github.com/tomsandovalmed/SyA-hub.git
cd SyA-hub
\end{lstlisting}

\subsection{Instalar dependencias}
\begin{lstlisting}[language=bash]
npm install
\end{lstlisting}

Dependencias de ejecución (\texttt{dependencies}):
\begin{center}
\rowcolors{2}{salight}{white}
\begin{tabularx}{\linewidth}{lX}
\toprule
\textbf{Paquete} & \textbf{Uso} \\
\midrule
\texttt{express}               & Servidor HTTP y API REST \\
\texttt{vite}                  & Servidor de desarrollo (middleware) y compilación \\
\texttt{@vitejs/plugin-react}  & Soporte de React/JSX en Vite \\
\texttt{@tailwindcss/vite}     & Integración de Tailwind CSS \\
\texttt{react}, \texttt{react-dom} & Librería de interfaz \\
\texttt{chart.js}, \texttt{react-chartjs-2} & Gráficos de líneas de los contadores \\
\texttt{lucide-react}          & Iconografía \\
\bottomrule
\end{tabularx}
\end{center}

Dependencias de desarrollo (\texttt{devDependencies}): \texttt{typescript}, \texttt{tsx}
(ejecuta \ruta{server.ts} sin compilar), \texttt{esbuild}, \texttt{tailwindcss},
\texttt{vite-plugin-javascript-obfuscator} y los tipos \texttt{@types/*}.

\subsection{Scripts disponibles}
\begin{center}
\rowcolors{2}{salight}{white}
\begin{tabularx}{\linewidth}{lX}
\toprule
\textbf{Comando} & \textbf{Qué hace} \\
\midrule
\texttt{npm run dev}   & Inicia el servidor en modo desarrollo con recarga en caliente \\
\texttt{npm run lint}  & Verifica tipos con \texttt{tsc --noEmit} \\
\texttt{npm run build} & Compila y ofusca el frontend y empaqueta el servidor en \ruta{dist/} \\
\texttt{npm run start} & Ejecuta la versión compilada (\ruta{dist/server.js}) \\
\texttt{npm run clean} & Elimina la carpeta \ruta{dist/} \\
\bottomrule
\end{tabularx}
\end{center}

\subsection{Modo desarrollo}
\begin{lstlisting}[language=bash]
npm run dev
\end{lstlisting}
Abra \url{http://localhost:3000}. El puerto está fijado en la constante
\texttt{PORT} de \ruta{server.ts}; cámbielo allí si lo necesita.

\subsection{Modo producción}
\begin{lstlisting}[language=bash]
npm run build
npm run start
\end{lstlisting}

\begin{aviso}
El servidor decide el modo así: si \texttt{NODE\_ENV=production} \textbf{o si existe la
carpeta \ruta{dist/}}, sirve la versión compilada. Por eso, después de un
\texttt{npm run build}, el comando \texttt{npm run dev} seguirá mostrando la versión
compilada. Ejecute \texttt{npm run clean} para volver al modo desarrollo.
\end{aviso}

\subsection{Exposición pública temporal (opcional)}
Para mostrar la aplicación fuera de la red local puede usarse un túnel:
\begin{lstlisting}[language=bash]
cloudflared tunnel --url http://localhost:3000
npx localtunnel --port 3000
\end{lstlisting}
Recuerde que las credenciales actuales son de demostración (sección~\ref{sec:credenciales}):
no exponga la aplicación públicamente sin cambiarlas.

% ============================================================================
\section{Guía de uso}

\subsection{Inicio de sesión}
Al abrir la aplicación se muestra la pantalla \emph{Iniciar Sesión}. Tras autenticarse,
el usuario es dirigido según su rol (ver sección~\ref{sec:rbac}). La sesión vive en
memoria del navegador: al recargar la página se vuelve a pedir el inicio de sesión.

\subsection{Centro de Herramientas (Hub)}
Pantalla principal para \texttt{ADMIN} y \texttt{WORKER}. La barra superior da acceso a
\emph{Herramientas}, \emph{Historial}, \emph{Tareas} y, solo para administradores,
\emph{Admin}. El menú de perfil (esquina superior derecha) muestra nombre, correo, rol
y el botón \emph{Cerrar sesión}.

\subsection{Intérprete Perfmon SQL}
\begin{enumerate}[leftmargin=*]
  \item En el Hub pulse \textbf{Utilizar Herramienta} y elija \textbf{Módulo SII} o
        \textbf{Módulo TGR}.
  \item Pulse \textbf{Cargar archivo(s) o carpeta} y seleccione uno o varios CSV de Perfmon.
  \item Pulse \textbf{Procesar archivo(s)}. Si se cargan varios archivos, los contadores
        con el mismo identificador se combinan en una sola serie.
  \item Navegue entre contadores con las pestañas; cada uno muestra un gráfico, una tabla
        con mínimo, promedio, máximo y desviación estándar, y la lista de alertas.
        \emph{Ver en gráfico} resalta el punto exacto de la alerta.
  \item Use el selector \textbf{Servidor} para filtrar si el log contiene varios equipos.
\end{enumerate}

\paragraph{Contadores reconocidos.} El analizador busca en los encabezados del CSV estas
cadenas (archivo \ruta{src/modules/analizador-rendimiento/utils/procesadorCsv.ts},
constante \texttt{TARGET\_METRICS}):
\texttt{Memory\textbackslash Available MBytes}, \texttt{\% Processor Time},
\texttt{Full Scans/sec}, \texttt{Buffer cache hit ratio}, \texttt{Transactions/sec} y
\texttt{User Connections}. Para agregar más contadores, añada entradas a esa lista.

\begin{aviso}[Archivos .blg]
El selector acepta \texttt{.blg}, pero el analizador lee el archivo como texto CSV.
Los \texttt{.blg} son binarios, por lo que deben convertirse a CSV antes de cargarlos.
En Windows suele hacerse con la herramienta \texttt{relog}
(\texttt{relog archivo.blg -f csv -o archivo.csv}); verifique la sintaxis en la
documentación de Microsoft de su versión de Windows.
\end{aviso}

\subsection{Umbrales y alertas}
El botón \textbf{Configuración} abre los umbrales de análisis. Al guardar, los
contadores ya cargados se reevalúan inmediatamente. Valores por defecto (definidos en
\texttt{initialThresholds}, archivo \ruta{src/modules/analizador-rendimiento/data/datosSimulados.ts}):

\begin{center}
\rowcolors{2}{salight}{white}
\begin{tabular}{lcc}
\toprule
\textbf{Métrica} & \textbf{Aviso} & \textbf{Crítico} \\
\midrule
\% Processor Time            & $\geq 50\%$ & $\geq 80\%$ \\
Buffer cache hit ratio       & $\leq 95\%$ & $\leq 80\%$ \\
Memory Pages/sec (límite)    & \multicolumn{2}{c}{20} \\
Cache Hit Ratio OLTP / OLAP  & \multicolumn{2}{c}{95\% / 80\%} \\
SQL Compilations/sec (límite)& \multicolumn{2}{c}{100} \\
Lock Requests/sec (límite)   & \multicolumn{2}{c}{1000} \\
\bottomrule
\end{tabular}
\end{center}

\begin{nota}
En la versión actual, el detector de alertas
(\ruta{utils/detectorUmbrales.ts}) solo genera alertas para \textbf{CPU} y
\textbf{Buffer cache hit ratio}. Los demás umbrales se guardan y se envían a la API,
pero aún no disparan alertas: son el punto de extensión natural para nuevas reglas.
\end{nota}

\subsection{Reportes y exportación}
\begin{itemize}[leftmargin=*]
  \item \textbf{Generar Reporte Completo}: descarga un CSV (para Excel) o un JSON (para
        integración con otros sistemas) con contadores, estadísticas y alertas.
  \item \textbf{Descargar gráfico}: exporta el gráfico activo como PNG, con el nombre
        del servidor y del contador en la cabecera.
  \item \textbf{Más acciones}: copia al portapapeles la condición o las estadísticas.
\end{itemize}

\subsection{Repositorio de archivos cargados}
Cada procesamiento queda guardado en el \texttt{localStorage} del navegador. Si el
análisis es demasiado grande para el almacenamiento local, solo se conservan los
metadatos y será necesario volver a procesar el archivo para ver sus datos.

\subsection{Historial y Notas/Tareas}
El \emph{Historial} registra cada análisis realizado. \emph{Tareas} permite crear notas
con título, descripción y estado. Ambos se mantienen en memoria y se reinician al
recargar la página.

\subsection{Portal de clientes}
Los usuarios con rol \texttt{CLIENT} ven exclusivamente el \emph{Centro de informes
corporativo} con los informes cuyo \texttt{clientId} coincide con el suyo. Pueden buscar
por texto, filtrar por tipo y descargar cada informe.

\subsection{Panel de administración}
Disponible solo para \texttt{ADMIN}. Permite crear usuarios (nombre, correo, rol,
\texttt{clientId} opcional y contraseña opcional), editarlos, cambiar el rol directamente
desde la tabla y eliminarlos (con confirmación).

\begin{aviso}
Los cambios del panel de administración se guardan \textbf{solo en la memoria del
navegador} y se pierden al recargar la página. Además, no llegan al servidor: un usuario
creado desde el panel inicia sesión gracias a la validación local de respaldo del
frontend (ver sección~\ref{sec:seguridad}).
\end{aviso}

% ============================================================================
\section{Control de acceso por roles (RBAC)}\label{sec:rbac}

El sistema define tres roles en el tipo \texttt{Role} de \ruta{src/types.ts}:

\begin{lstlisting}
export type Role = 'ADMIN' | 'WORKER' | 'CLIENT';
\end{lstlisting}

\subsection{Matriz de permisos}
\begin{center}
\rowcolors{2}{salight}{white}
\begin{tabularx}{\linewidth}{Xccc}
\toprule
\textbf{Funcionalidad} & \textbf{ADMIN} & \textbf{WORKER} & \textbf{CLIENT} \\
\midrule
Pantalla inicial tras el login      & Panel Admin & Hub & Portal cliente \\
Centro de Herramientas (Hub)        & Sí & Sí & No \\
Intérprete Perfmon SII / TGR        & Sí & Sí & No \\
Configurar umbrales y exportar      & Sí & Sí & No \\
Historial de acciones               & Sí & Sí & No \\
Notas y tareas                      & Sí & Sí & No \\
Panel de administración de usuarios & Sí & No & No \\
Portal de informes del cliente      & No & No & Sí (solo su \texttt{clientId}) \\
\bottomrule
\end{tabularx}
\end{center}

\subsection{Dónde se aplica}
\begin{itemize}[leftmargin=*]
  \item \ruta{src/App.tsx}: tras el login, \texttt{handleLoginSuccess} envía a cada rol a su
        pantalla inicial. Si el rol es \texttt{CLIENT} se renderiza únicamente
        \texttt{ClientReportsView}; el resto de la aplicación no se monta. La vista
        \texttt{admin} solo se renderiza si \texttt{authUser.role === "ADMIN"}.
  \item \ruta{src/components/ui/TopNavBar.tsx}: el botón \emph{Admin} solo se muestra
        al rol \texttt{ADMIN}.
  \item \ruta{src/core/auth.ts} y \ruta{server.ts}: \texttt{getClientReports} y el endpoint
        \texttt{GET /api/v1/reports?clientId=...} filtran los informes por \texttt{clientId}.
\end{itemize}

\subsection{Cómo modificar los roles}
\begin{enumerate}[leftmargin=*]
  \item \textbf{Agregar un rol} (por ejemplo \texttt{AUDITOR}): añádalo al tipo
        \texttt{Role} en \ruta{src/types.ts}.
  \item Defina su pantalla inicial en \texttt{handleLoginSuccess} de \ruta{src/App.tsx} y
        las condiciones de renderizado de las vistas que podrá ver.
  \item Muestre u oculte los botones de navegación en \ruta{TopNavBar.tsx}.
  \item Agregue la opción en los \texttt{<select>} de rol de
        \ruta{src/views/AdminDashboardView.tsx}.
  \item Ejecute \texttt{npm run lint}: TypeScript señalará los lugares que falten.
\end{enumerate}

% ============================================================================
\section{Credenciales actuales}\label{sec:credenciales}

\subsection{Cuentas de demostración}
\begin{center}
\rowcolors{2}{salight}{white}
\begin{tabular}{llll}
\toprule
\textbf{Rol} & \textbf{Correo} & \textbf{Contraseña} & \textbf{clientId} \\
\midrule
ADMIN  & \texttt{admin@sachile.cl}       & \texttt{admin123}  & --- \\
WORKER & \texttt{worker@sachile.cl}      & \texttt{worker123} & --- \\
CLIENT & \texttt{cliente.sii@sachile.cl} & \texttt{client123} & \texttt{SII} \\
CLIENT & \texttt{cliente.tgr@sachile.cl} & \texttt{client123} & \texttt{TGR} \\
\bottomrule
\end{tabular}
\end{center}

\subsection{Dónde están alojadas}
Todas las credenciales viven en el \textbf{código fuente}; no hay base de datos ni archivo
\texttt{.env} involucrado.

\begin{center}
\rowcolors{2}{salight}{white}
\begin{tabularx}{\linewidth}{lX}
\toprule
\textbf{Archivo} & \textbf{Contenido} \\
\midrule
\ruta{src/config/auth.ts} &
  Arreglo \texttt{initialUsers}: id, correo, nombre, rol y \texttt{clientId} de cada
  usuario. También \texttt{demoReports}, los informes visibles en el portal de clientes. \\
\ruta{server.ts} (endpoint \texttt{/api/v1/auth/login}) &
  Regla de contraseña del servidor. \\
\ruta{src/core/auth.ts} (\texttt{authenticateUser}) &
  La misma regla, usada como respaldo en el navegador. \\
\ruta{src/views/LoginView.tsx} (\texttt{passwordHint}) &
  Texto de ayuda que muestra la contraseña esperada según el correo escrito. \\
\ruta{src/views/AdminDashboardView.tsx} (\texttt{getCredentialHint}) &
  Columna \emph{Credencial demo} de la tabla de usuarios. \\
\bottomrule
\end{tabularx}
\end{center}

\paragraph{Regla de contraseña.} Los usuarios iniciales no tienen contraseña propia; se
usa una contraseña derivada del rol:
\begin{lstlisting}
const expected = user.password ?? `${user.role.toLowerCase()}123`;
\end{lstlisting}
Es decir, \texttt{admin123}, \texttt{worker123} o \texttt{client123}. Si un usuario tiene
el campo \texttt{password}, se usa ese valor.

\subsection{Cómo cambiarlas}
\begin{enumerate}[leftmargin=*]
  \item \textbf{Cambiar correos, nombres o agregar usuarios}: edite \texttt{initialUsers}
        en \ruta{src/config/auth.ts}.
  \item \textbf{Asignar una contraseña propia}: agregue el campo \texttt{password} al
        usuario, por ejemplo:
\begin{lstlisting}
{
  id: 'admin-1',
  email: 'admin@sachile.cl',
  name: 'Administrador S&A',
  role: 'ADMIN',
  password: 'NuevaClaveSegura!2026',
},
\end{lstlisting}
  \item \textbf{Cambiar la regla por rol}: modifique la línea \texttt{expected} en
        \ruta{server.ts} \emph{y} en \ruta{src/core/auth.ts} (deben coincidir).
  \item \textbf{Quitar las pistas de contraseña}: elimine o reemplace
        \texttt{passwordHint} en \ruta{LoginView.tsx} y la columna \emph{Credencial demo}
        en \ruta{AdminDashboardView.tsx}.
  \item Reinicie el servidor (\texttt{npm run dev}) o vuelva a compilar
        (\texttt{npm run build}) para que los cambios surtan efecto.
\end{enumerate}

% ============================================================================
\section{Consideraciones de seguridad}\label{sec:seguridad}

La autenticación actual es una \textbf{simulación para demostración}. Antes de usar la
aplicación con datos reales o exponerla en Internet considere lo siguiente:

\begin{itemize}[leftmargin=*]
  \item \textbf{Las credenciales viajan en el frontend.} \ruta{src/config/auth.ts} se
        incluye en el JavaScript que descarga el navegador. La ofuscación dificulta
        leerlo, pero no lo impide.
  \item \textbf{Validación de respaldo en el cliente.} Si el servidor rechaza el login,
        \texttt{authenticateUser} vuelve a validar localmente en el navegador; por lo
        tanto la decisión final de acceso ocurre en el cliente.
  \item \textbf{Token simulado.} El servidor responde con un token del tipo
        \texttt{mock-<id>} que no se verifica en ningún endpoint; la API no exige
        autenticación.
  \item \textbf{Las pistas de contraseña} están visibles en la pantalla de login.
  \item \textbf{Sin persistencia.} Usuarios, historial y tareas se reinician al recargar
        la página o reiniciar el servidor.
\end{itemize}

\begin{nota}[Siguiente paso sugerido]
El propio código deja preparado el camino para conectar un proveedor de autenticación
real y una base de datos (las respuestas de la API mencionan Firebase Auth y
PostgreSQL). Lo recomendable es: mover usuarios y contraseñas (con \emph{hash}) a una
base de datos, eliminar la validación de respaldo del frontend, emitir tokens firmados y
verificar el rol en cada endpoint del servidor.
\end{nota}

% ============================================================================
\section{API REST del servidor}

Todas las rutas están en \ruta{server.ts} y responden JSON.

\begin{center}
\rowcolors{2}{salight}{white}
\begin{tabularx}{\linewidth}{llX}
\toprule
\textbf{Método} & \textbf{Ruta} & \textbf{Descripción} \\
\midrule
GET  & \texttt{/api/v1/health}            & Estado del servicio \\
POST & \texttt{/api/v1/auth/login}        & Login con \texttt{\{email, password\}} \\
GET  & \texttt{/api/v1/reports}           & Informes; filtro opcional \texttt{?clientId=} \\
GET  & \texttt{/api/v1/perfmon/counters}  & Contadores de ejemplo; \texttt{?module=SII|TGR} \\
GET  & \texttt{/api/v1/perfmon/alerts}    & Alertas de ejemplo; \texttt{?module=SII|TGR} \\
GET  & \texttt{/api/v1/perfmon/history}   & Historial inicial \\
POST & \texttt{/api/v1/perfmon/analyze}   & Análisis simulado (datos de ejemplo) \\
POST & \texttt{/api/v1/perfmon/thresholds}& Recibe la configuración de umbrales (sin persistir) \\
\bottomrule
\end{tabularx}
\end{center}

\begin{nota}
El análisis real de los CSV ocurre en el navegador
(\ruta{procesadorCsv.ts}); los endpoints \texttt{perfmon/*} devuelven datos simulados de
\ruta{datosSimulados.ts} y sirven como contrato para un futuro backend.
\end{nota}

% ============================================================================
\section{Estructura del proyecto y puntos de personalización}\label{sec:estructura}

\begin{lstlisting}
SyA-hub/
|-- server.ts                 Servidor Express + API REST
|-- index.html                HTML base (titulo y favicon)
|-- vite.config.ts            Configuracion de Vite y ofuscacion
|-- public/S&A.svg            Favicon / logotipo
|-- docs/                     Este manual
`-- src/
    |-- App.tsx               Navegacion principal y reglas de rol
    |-- types.ts              Tipos User, Role, ReportItem
    |-- config/auth.ts        Usuarios e informes de demostracion
    |-- core/auth.ts          Login y obtencion de informes (cliente)
    |-- components/ui/        Barra superior y pie de pagina
    |-- views/                Login, Hub, Historial, Tareas, Admin, Portal cliente
    `-- modules/analizador-rendimiento/
        |-- VistaAnalizadorRendimiento.tsx
        |-- components/       Grafico, alertas, modales, seleccion de modulo
        |-- hooks/            Logica de carga, repositorio local y estado
        |-- utils/            Parser CSV y detector de umbrales
        |-- data/             Datos simulados y umbrales por defecto
        `-- types/            Tipos del analizador
\end{lstlisting}

\begin{center}
\rowcolors{2}{salight}{white}
\begin{tabularx}{\linewidth}{Xl}
\toprule
\textbf{Quiero cambiar\ldots} & \textbf{Archivo} \\
\midrule
Usuarios, roles iniciales e informes de clientes & \ruta{src/config/auth.ts} \\
Regla de contraseñas            & \ruta{server.ts} y \ruta{src/core/auth.ts} \\
Puerto del servidor             & \ruta{server.ts} (\texttt{PORT}) \\
Umbrales por defecto            & \ruta{.../data/datosSimulados.ts} \\
Reglas de alerta                & \ruta{.../utils/detectorUmbrales.ts} \\
Contadores Perfmon reconocidos  & \ruta{.../utils/procesadorCsv.ts} \\
Textos del Hub / tarjetas       & \ruta{src/views/InicioVista.tsx} \\
Descripción de módulos SII/TGR  & \ruta{.../components/SeleccionModulo.tsx} \\
Colores corporativos            & Clases Tailwind (p.\,ej. \texttt{\#002395}) en cada vista \\
Título de pestaña y favicon     & \ruta{index.html}, \ruta{public/S\&A.svg} \\
Nivel de ofuscación del build   & \ruta{vite.config.ts} \\
\bottomrule
\end{tabularx}
\end{center}

% ============================================================================
\section{Solución de problemas}

\begin{description}[leftmargin=*,style=nextline]
  \item[El puerto 3000 está ocupado] Cierre el otro proceso o cambie \texttt{PORT} en
        \ruta{server.ts}.
  \item[\texttt{npm run dev} muestra una versión antigua] Existe \ruta{dist/}; ejecute
        \texttt{npm run clean}.
  \item[``El archivo CSV no contiene contadores válidos''] El CSV no incluye ninguno de
        los contadores reconocidos, o es un \texttt{.blg} sin convertir.
  \item[``Contraseña incorrecta''] Revise la tabla de la sección~\ref{sec:credenciales};
        si cambió la regla, verifique que \ruta{server.ts} y \ruta{src/core/auth.ts}
        coincidan.
  \item[Los usuarios creados desaparecieron] Es el comportamiento esperado: se guardan
        solo en memoria (ver sección~\ref{sec:seguridad}).
\end{description}

\end{document}
